\section{Paying the Cost}

\begin{quotex}
It takes years of training to teach us to deal intelligently with the world of everyday life. Our schooling --- whether in plain reasoning or formal topics --- is rigorous, because the knowledge we are trying to impart is very complex. The same criteria apply to the sorcerers' world: their schooling, which relies on oral instruction and the manipulation of awareness, although different from ours, is just as rigorous, because their knowledge is as, or perhaps more, complex.
\flright{\textsc{Carlos Castenada}, \emph{The Power of Silence}}
\end{quotex}

Most people do not spend time mastering some field of knowledge, because they remain content to repeat whatever they have learned over and over, year after year. Yet when it comes to the most abstruse topics in metaphysics or philosophy, many people assume it is possible to become an instant expert.

Au contraire, it takes years of study and practice to begin to master metaphysics.

\paragraph{Gladwell's Law}


\textbf{Malcolm Gladwell}, in his blockbuster book \textit{Outliers}, proposed the rule following rule for success:

\textbf{It takes 10,000 hours of intensive practice to achieve mastery of complex skills and materials.}

That is equivalent to 5 years of 40 hour workweeks, but that is approximately what it takes to develop world class expertise of a talent. Obviously, after work and family obligations, there is little spare time for a hobby or other avocation. That is why we like to hear stories of novelists, for example, who had to wrote at odd times of day, while not engaged in family and/or work duties. In times past, mathematical and scientific discoveries were often made by men who were otherwise engaged.

This applies a fortiori to one's spiritual life, particularly the path of the Spiritual Knight and Hermetic philosophy. Few are able to dedicate adequate time.

Since Gladwell proposed his law, there have been emendations and criticisms. However, they are the exceptions that prove the rule. There are two, in particular, that have relevance.

\textbf{Success}. Talent is not sufficient to guarantee success, at least in worldly terms. Sometimes someone is aided by a wealthy benefactor, or a talent agent, or even a serendipitous break. A good friend of mine has been enjoying great guitarists in concert for years. So I recently asked him who are the great guitarists of today, and this was his response:

\begin{quotex}
Music biz is $>$ 90\% promoter agent vs 10\% performer. There are outstanding musicians out there. Too few care to look for them. So there is no buzz. Within 8 days, 9/14- 9/22, I saw live Kurt Rosenwinkel, DiMeola and Richard Thompson. Times change How many great living violinists can you name?
\end{quotex}

Although I listen to a lot of classical music, I cannot name the best living violinists; perhaps that young girl who is a street performer. Moreover, I cannot name one great oboist, ever in history. But I don't think it takes 10,000 hours to master the oboe.

\textbf{Teacher}. A good teacher or coach can save many of those hours.

It is of no value to repeat the same thing for all those hours, otherwise mailmen would get really good at mail delivery. We used to have the same mailman for 17 years. I suspect that he learned his route within the first few months.

So a good teacher can evaluate performance, make corrections, suggest alternatives, and so on, which might save years of effort if you had to do it on your own.

\paragraph{Sacred Science}

There is one other factor, without an easy explanation and which may offend our contemporary notions of egalitarianism, is that not everyone has the potential to excel in every field. In other words, we are not born a ``blank slate'' as John Locke proposed. Instead we have innate potentialities which may or may not be developed.

This is most especially true when it comes to metaphysics, or Sacred Science. Few have the aptitude for it. Many try, but then drift away. They learn the vocabulary and some basic ideas, but slowly drift away. For them, it was never more than a pastime. On the other hand, a few become very interested and spend as much time as they can. Rene Guenon, in Man and his Becoming, explains the way it is:

\begin{quotex}
Among those who receive the same teaching, each one understands and assimilates it more or less completely, more or less deeply, according to the extent of his own intellectual possibilities. This is how the selection quite naturally operates, without which there can be no true hierarchy.
\end{quotex}


\paragraph{The Sun and the Moon}

One of the stumbling blocks to the study of metaphysics is that, at the superficial level, it ``looks like'' philosophy. That is, it seems to be a subject matter that can be debated and discussed, just like philosophical, political, or scientific theories. At one level, that is correct. But that type of discursive knowledge is just indirect knowledge.

Profane teachings to not require a change in the learner. That is, there is no necessity for moral purification to become a great physicist. However, that is not the case for the study of metaphysics. The physicist observed the outer world and performs experiments.

Analogously, for the metaphysician carefully observes, not the outer world, but the inner world of sensations, emotions, images, and thoughts. By accumulating data of this type, he learns the concepts, not indirectly, but through direct observations, called ``intuition''. Without have made the requisite efforts of this type, the metaphysical teachings can never be fully known. Guenon explains:

\begin{quotex}
in the transcendent order, only direct perception or inspiration attains the principle itself, i.e., what is highest, and all that remains is to draw out the consequences and the various applications.
\end{quotex}


This is represented by symbolism:

\begin{quotex}
The intuitive faculty and the discursive faculty, may be represented respectively by the sun and the moon.
\end{quotex}


Discursive knowledge is always open to debate; intuitive knowledge is certain. The Shaman explains:

\begin{quotex}
Don Juan had stated his belief that the Christian idea of being cast out from the Garden of Eden sounded to him like an allegory for losing our silent knowledge, our knowledge of intent. Sorcery, then, was a going back to the beginning, a return to paradise.

``What we need to do to allow magic to get hold of us is to banish doubt from our minds,'' he said. ``Once doubts are banished, anything is possible.''
\end{quotex}


Who has 10 thousand hours for that result?

\url{https://www.youtube.com/watch?v=oUTwVtS\_-NA}

\flrightit{Posted on 2022-09-25 by Cologero}

\begin{center}* * *\end{center}

\begin{footnotesize}
\begin{sffamily}
\texttt{Tannheuser on 2022-09-26 at 10:54 said:}

A good teacher can save his students thousands of hours. The importance of having a teacher is that instead of just dispensing general advice (which can be found in books or online), the teacher gets to know the student better than they know themselves, and can provide invaluable personalized guidance.

One example is the way that good teachers often don't answer students' questions directly, but instead guide the students toward finding the answers themselves. In many cases, the work necessary to discover the answer is as important as the answer itself.

From observing those more advanced than myself, I believe that to progress past a certain advanced stage in esoteric development, it is necessary to become a teacher oneself.

\hfill

\texttt{mark siegmund on 2022-09-26 at 16:50 said:}

'' A Case Western Reserve University team has subsequently performed a review of over 9,000 research papers about practice relating to acquiring skills. In their paper, they found that deliberate practice explained 26\% of the variance in performance for games, 21\% for music, 18\% for sports, 4\% for education, and less than 1\% for professions. They conclude that deliberate practice is important, but not as important as has been argued in Gladwell's book.''


\hfill

\end{sffamily}
\end{footnotesize}

